% generated by src/make_tables.py -- do not edit
\begin{table}[!t]
\centering
\caption{Komposisi dan struktur himpunan data. $d(G)$ adalah kepadatan teramati, $\bar{C}$ koefisien klaster rata-rata.}
\label{tab:dataset}
\footnotesize
\setlength{\tabcolsep}{4pt}
\begin{tabular}{@{}lrrrrr@{}}
\toprule
Blok & $n$ & $|V|$ & $|E|$ & $d(G)$ & $\bar{C}$ \\
\midrule
A (acak) & 600 & 20--100 & 10--3533 & 0.053--0.789 & 0.397 \\
B, kecil & 30 & 24 & 84--87 & 0.304--0.315 & 0.759 \\
B, menengah & 30 & 40 & 303--368 & 0.388--0.472 & 0.771 \\
B, besar & 30 & 60 & 845--1029 & 0.477--0.581 & 0.797 \\
C (nyata) & 1 & 156 & 1029 & 0.085 & 0.649 \\
\bottomrule
\end{tabular}
\end{table}
